% OpenStax Calculus Volume 2 -- Section 2.2 Exercises: Determining Volumes by Slicing
% Exercises reproduced from OpenStax, licensed under CC BY 4.0.
% Source: https://openstax.org/books/calculus-volume-2/pages/2-2-determining-volumes-by-slicing
\documentclass{exercises}
\title{OpenStax Calculus Volume 2}
\subtitle{Section 2.2 Exercises: Determining Volumes by Slicing}
\sourceurl{https://openstax.org/books/calculus-volume-2/pages/2-2-determining-volumes-by-slicing}
\author{}
\date{}

\begin{document}
\maketitle

\begin{enumerate}[start=1]
\item Derive the formula for the volume of a sphere using the slicing method.
\item Use the slicing method to derive the formula for the volume of a cone.
\item Use the slicing method to derive the formula for the volume of a tetrahedron with side length $a$.
\item Use the disk method to derive the formula for the volume of a trapezoidal cylinder.
\item Explain when you would use the disk method versus the washer method. When are they interchangeable?
\end{enumerate}

For the following exercises, draw a typical slice and find the volume using the slicing method for the given volume.

\begin{enumerate}[resume]
\item A pyramid with height 6 units and square base of side 2 units, as pictured here. \\ \textit{[Figure: This figure is a pyramid with base width of 2 and height of 6 units.]}
\item A pyramid with height 4 units and a rectangular base with length 2 units and width 3 units, as pictured here. \\ \textit{[Figure: This figure is a pyramid with base width of 2, length of 3, and height of 4 units.]}
\item A tetrahedron with a base side of 4 units, as seen here. \\ \textit{[Figure: This figure is an equilateral triangle with side length of 4 units.]}
\item A pyramid with height 5 units, and an isosceles triangular base with lengths of 6 units and 8 units, as seen here. \\ \textit{[Figure: This figure is a pyramid with a triangular base. The view is of the base. The sides of the triangle measure 6 units, 8 units, and 8 units. The height of the pyramid is 5 units.]}
\item A cone of radius $r$ and height $h$ has a smaller cone of radius $r/2$ and height $h/2$ removed from the top, as seen here. The resulting solid is called a frustum. \\ \textit{[Figure: This figure is a 3-dimensional graph of an upside down cone. The cone is inside of a rectangular prism that represents the $xyz$ coordinate system. The radius of the bottom of the cone is ``r'' and the radius of the top of the cone is labeled ``r/2''.]}
\end{enumerate}

For the following exercises, draw an outline of the solid and find the volume using the slicing method.

\begin{enumerate}[resume]
\item The base is a circle of radius $a$. The slices perpendicular to the base are squares.
\item The base is a triangle with vertices $(0,0),(1,0)$, and $(0,1)$. Slices perpendicular to the $x$-axis are semicircles.
\item The base is the region under the parabola $y=1-x^{2}$ in the first quadrant. Slices perpendicular to the xy-plane and parallel to the $y$-axis are squares.
\item The base is the region under the parabola $y=1-x^{2}$ and above the $x$-axis. Slices perpendicular to the $y$-axis are squares.
\item The base is the region enclosed by $y=x^{2}$ and $y=9$. Slices perpendicular to the $x$-axis are right isosceles triangles. The intersection of one of these slices and the base is the leg of the triangle.
\item The base is the area between $y=x$ and $y=x^{2}$. Slices perpendicular to the $x$-axis are semicircles.
\end{enumerate}

For the following exercises, draw the region bounded by the curves. Then, use the disk method to find the volume when the region is rotated around the $x$-axis.

\begin{enumerate}[resume]
\item $x+y=8,x=0,\,\text{and}\,y=0$
\item $y=2x^{2},x=0,x=4,\,\text{and}\,y=0$
\item $y=e^{x}+1,x=0,x=1,\,\text{and}\,y=0$
\item $y=x^{4},x=0,\,\text{and}\,y=1\,\text{for}\,x\ge0$
\item $y=\sqrt{x},x=0,x=4,\,\text{and}\,y=0$
\item $y=\sin x,y=\cos x,\,\text{and}\,x=0$
\item $y=\frac{1}{x},x=2,\,\text{and}\,y=3$
\item $x^{2}-y^{2}=9\,\text{and}\,x+y=9,y=0\,\text{and}\,x=0$
\end{enumerate}

For the following exercises, draw the region bounded by the curves. Then, find the volume when the region is rotated around the $y$-axis.

\begin{enumerate}[resume]
\item $y=4-\frac{1}{2}x,x=0,\,\text{and}\,y=0$
\item $y=2x^{3},x=0,x=1,\,\text{and}\,y=0$
\item $y=3x^{2},x=0,\,\text{and}\,y=3$
\item $y=\sqrt{4-x^{2}},y=0,\,\text{and}\,x=0$
\item $y=\frac{1}{\sqrt{x+1}},x=0,\,x=3,\,\text{and}\,y=0$
\item $x=\sec(y)\,\text{and}\,y=\frac{\pi}{4},\,y=0\,\text{and}\,x=0$
\item $y=\frac{1}{x+1},x=0,\,x=2,\,\text{and}\,y=0$
\item $y=4-x,y=x,\,\text{and}\,x=0$
\end{enumerate}

For the following exercises, draw the region bounded by the curves. Then, find the volume when the region is rotated around the $x$-axis.

\begin{enumerate}[resume]
\item $y=x+2,y=x+6,x=0,\,\text{and}\,x=5$
\item $y=x^{2}\,\text{and}\,y=x+2$
\item $x^{2}=y^{3}\,\text{and}\,x^{3}=y^{2}$
\item $y=4-x^{2}\,\text{and}\,y=2-x$
\item \techrequired $y=\cos x,y=e^{-x},x=0,\,\text{and}\,x=1.2927$
\item $y=\sqrt{x}\,\text{and}\,y=x^{2}$
\item $y=\sin x,\,y=5\sin x,x=0\,\text{and}\,x=\pi$
\item $y=\sqrt{1+x^{2}}\,\text{and}\,y=\sqrt{4-x^{2}}$
\end{enumerate}

For the following exercises, draw the region bounded by the curves. Then, use the washer method to find the volume when the region is revolved around the $y$-axis.

\begin{enumerate}[resume]
\item $y=\sqrt{x},x=4,\,\text{and}\,y=0$
\item $y=x+2,y=2x-1,\,\text{and}\,x=0$
\item $y=\sqrt[3]{x}\,\text{and}\,y=x^{3}$
\item $x=e^{2y},x=y^{2},y=0,\,\text{and}\,y=\ln(2)$
\item $x=\sqrt{9-y^{2}},x=e^{-y},y=0,\,\text{and}\,y=3$
\item Yogurt containers can be shaped like frustums. Rotate the line $y=\frac{1}{m}x$ around the $y$-axis to find the volume between $y=a\,\text{and}\,y=b$. \\ \textit{[Figure: This figure has two parts. The first part is a solid cone. The base of the cone is wider than the top. It is shown in a 3-dimensional box. Underneath the cone is an image of a yogurt container with the same shape as the figure.]}
\item Rotate the ellipse $(x^{2}/a^{2})+(y^{2}/b^{2})=1$ around the $x$-axis to approximate the volume of a football, as seen here. \\ \textit{[Figure: This figure has an oval that is approximately equal to the image of a football.]}
\item Rotate the ellipse $(x^{2}/a^{2})+(y^{2}/b^{2})=1$ around the $y$-axis to approximate the volume of a football.
\item A better approximation of the volume of a football is given by the solid that comes from rotating $y=\sin x$ around the $x$-axis from $x=0$ to $x=\pi$. What is the volume of this football approximation, as seen here? \\ \textit{[Figure: This figure has a 3-dimensional oval shape. It is inside of a box parallel to the $x$-axis on the bottom front edge of the box. The $y$-axis is vertical to the solid.]}
\item What is the volume of the Bundt cake that comes from rotating $y=\sin x$ around the $y$-axis from $x=0$ to $x=\pi$? \\ \textit{[Figure: This figure is a graph of a 3-dimensional solid. It is round, bigger towards the bottom. It has a hole in the center that progressively gets smaller towards the bottom. Next to the graph is an image of a bundt cake, resembling the solid.]}
\end{enumerate}

For the following exercises, find the volume of the solid described.

\begin{enumerate}[resume]
\item The base is the region between $y=x$ and $y=x^{2}$. Slices perpendicular to the $x$-axis are semicircles.
\item The base is the region enclosed by the generic ellipse $(x^{2}/a^{2})+(y^{2}/b^{2})=1$. Slices perpendicular to the $x$-axis are semicircles.
\item Bore a hole of radius $a$ down the axis of a right cone of height $b$ and radius $b$ through the base of the cone as seen here. \\ \textit{[Figure: This figure is an upside down cone. It has a radius of the top as ``b'', center at ``a'', and height as ``b''.]}
\item Find the volume common to two spheres of radius $r$ with centers that are $2h$ apart, as shown here. \\ \textit{[Figure: This figure has two circles that intersect. Both circles have radius ``r''. There is a line segment from one center to the other. In the middle of the intersection of the circles is point ``h''. It is on the line segment.]}
\item Find the volume of a spherical cap of height $h$ and radius $r$ where $h<r$, as seen here. \\ \textit{[Figure: This figure a portion of a sphere. This spherical cap has radius ``r'' and height ``h''.]}
\item Find the volume of a sphere of radius $R$ with a cap of height $h$ removed from the top, as seen here. \\ \textit{[Figure: This figure is a sphere with a top portion removed. The radius of the sphere is ``R''. The distance from the center to where the top portion is removed is ``R-h''.]}
\end{enumerate}

\end{document}
